%% ma: L10-B14-0046
%% lop: 10
%% chuong: 5
%% bai: 14
%% dang: 10.14.10
%% loai: TL
%% muc: NB
%% dap_an: "-"
%% nguon: "chuyen-de-cac-so-dac-trung-cua-mau-so-lieu-khong-ghep-nhom-toan-10-chuong-trinh-moi.pdf (GV. Phan Nhật Linh) (P05), Bài 03 Dạng 2 Bài tập tự luận 6"
%% kien_thuc: "Lớp 10 Bài 13 (số trung bình), Bài 14 (phương sai, độ lệch chuẩn)"
%% trang_thai: chinh-thuc
%% ghi_chu: "Tự duyệt theo yêu cầu giáo viên 22/10/2026."
\begin{ex}
Theo kết quả thống kê điểm thi giữa kì 2 môn Toán khối 11 của một trường THPT, người ta tính được phương sai của mẫu số liệu là $s^2=0{,}573$. Tính độ lệch chuẩn của mẫu số liệu (làm tròn đến hàng phần nghìn).
\loigiai{$s=\sqrt{s^2}=\sqrt{0{,}573}\approx0{,}757$.}
\end{ex}
